\documentclass[11pt,a4paper,oneside]{article}
\usepackage{amsmath,amsthm,amssymb,amsfonts,adjustbox,bm}
%\usepackage[dvipsnames]{xcolor}
\usepackage{mathtools}
\usepackage{refstyle}
\usepackage[colorlinks=true, citecolor=blue, linkcolor=blue, urlcolor=blue,breaklinks]{hyperref}
\usepackage{pdflscape}
\usepackage[flushleft]{threeparttable}
\usepackage{multirow}
\usepackage{longtable}
\usepackage{eurosym}
\usepackage{enumerate}
\usepackage{booktabs}
\usepackage{siunitx}
\usepackage{rotating}
\usepackage{graphicx}
\usepackage{subcaption}
\usepackage[section]{placeins}
\usepackage{xcolor}
\usepackage{geometry}
\usepackage{changepage}
\usepackage[natbibapa]{apacite}


\title{Skills, Inequality, and Referrals}
\author{J. A. D. Contreras, M. M. Herrera, E. Reuben, R. Tuncer}%\footnote{}}

\begin{document}

\maketitle



% \section*{Abstract}


% \section{Introduction}



% \section{Literature Review}



\section{Experiment}

\subsection{Experimental Task}

The experiment consists of three interrelated parts, beginning with the individual skill measurement. In this part, each participant engages with two different skill tests in randomized order. Both skill tests have right and wrong answers and are limited to 5-minutes per test. Participants in this part are compensated for their performance, i.e., expected earnings weakly increase with the number of correct answers per test (see section \ref{Procedures} for details on procedures).

The first skill measure is a modified Raven's Standard Progressive Matrices test \citep{raven_performances_1936, raven_standard_1976}. We use the 18-item version also found in \citet{reuben_productivity_2024}. First half of the items have 6 possible answers while the remaining half have 8 such that difficulty of the test increases gradually. Raven's test is a well-established measure of fluid intelligence, a specific cognitive ability defining an individual's capacity to reason and solve problems in novel situations independent of past knowledge \citep{carpenter_what_1990, heckman_hard_2012, schilbach_psychological_2016}. 

% Performance in this test is likely to be productivity relevant.

The second skill measure is a novel, more race- and gender-inclusive version of the Reading the Mind in the Eyes Test (RMET) based on the original \citet{baron-cohen_reading_2001}. The test consists of cropped photos of faces, with only the eye-region visible. In each photo, a professional actor expresses an emotion. Participants need to choose the emotion that best describes the photo from four available options. We use the first 36 (out of 37) items of the Multiracial Reading the Mind in the Eyes (MRMET) developed by \citet{kim_multiracial_2022}.\footnote{We changed one out of the 4 possible emotions described in the practice image.} This test measures an individual's ability to recognize emotions in others \citep{oakley_theory_2016}, and has been previously used in economic experiments \citep{reuben_can_2009,weidmann_team_2021,van_leeuwen_predictably_2018}.  

After the skill measures, participants proceed with the second part of the experiment. In this part, participants are asked to refer three classmates for each skill test. Referrals can only be made among classmates, and we exclude self-referrals. Participants see a list with names of their classmates and tick the box in front of a name to refer an individual. They can freely refer up to six classmates except themselves, being able to refer the same three individuals twice. Participants are compensated based on their referrals' performance: Their expected earnings weakly increase by the number of top performers they include in their  referral choices. In other words, participants strategically refer among the high performers in their classroom. We experimentally manipulate the incentives to refer depending on the treatments, as explained in the section \ref{Treatments}.

% \textcolor{red}{This task is not important. Explain it briefly in the first section. it is also not varying by treatment so it does not fit here.}
In the third part of the experiment, participants engage with a final, performance-independent referral task where they make three additional referrals. The task in this part resembles the task in part two: Participants use the same list with the names of their classmates to refer classmates except themselves. But this time compensation depends on participants' own ability to identify low-SES individuals. Their expected earnings strictly increase by the number of low-SES classmates they include in their referral choices.   

Participants also fill a demographic questionnaire which collects information about their gender, parental education, and government-assigned stratum number (from 1 to 6) at the beginning of the experiment. The instructions, consent form, demographic questions, skill measures, and referral tasks can be found in the \textcolor{red}{Appendices 1-5}.   

\subsection{Treatments} \label{Treatments}

% \textcolor{red}{Explain we have two treatments. Explain that for the baseline we choose a top3 of the best classmates irrespective of other characteristics. Explain bonuses depend on whether the chosen referral is in the top 3. Then explain that the main difference between quotas and the baseline is that we have the top low class plus the top two others irrespective of other charactertistics. Also, bonuses depend on whether the chosen referral is in the top 3, but the composition of the top 3 is the main difference.}

In this experiment there are two exogenous treatments, introduced during the second part of the experiment where participants make performance-based referrals. Assignment to the treatments is based on participant's reported gender and stratum, collected at the beginning of the experiment. Participants in strata 1 to 3 (low SES) and 4-6 (high SES) are grouped together. In both treatments, participants are required to refer three individuals for each skill measure in randomized order. 

Participants in the \textbf{baseline} treatment are incentivized to refer individuals they predict to be among the \textit{top 3} for their classroom, independent of other characteristics. In the \textbf{quota} treatment, the incentives remains the same for the first two individuals among the \textit{top 3}, with one place reserved for the highest performing low-SES classmate. Participants are still incentivized to refer based on performance, with an additional incentive to  referring the best performer among the low-SES participants. The purpose of this treatment is to increase referrals among low-SES by giving participants an incentive to refer based on the SES of classmates. 




% Once participants across \textbf{baseline} and \textbf{quota} treatments made their referrals, they proceeded with a joint referral exercise. Participants in both treatment arms were asked to make three additional referrals. This time, the incentives to earn the bonus were not based on their referrals performance, but they depended on the participant's own ability to refer low-SES individuals: 10\% of participants in each classroom were randomly selected to take part in a lottery to win a bonus worth 100,000 Pesos (about \$26). Among those chosen for the lottery, we randomly selected one of the participant's three referrals, and if the selected referral was among the low-SES individuals, the participant received the bonus. The purpose of this exercise was to identify participants'  ability to recognize low-SES individuals in their classrooms.

% \subsection{Treatments}

% After participants completed the skill measures, they proceeded with the referral treatments. In each classroom, participants were randomly assigned to one of the two referral treatments based on their gender and SES. In both treatments, we asked participants to refer three individuals. Referrals could be made only among classmates, excluding referrals of their own selves. The incentivization scheme in both treatments were based on referrals' performance: For each skill measure, we ranked each participant's correct answers from the highest to the lowest number and designated the first 3 as the \textit{top three} performers for that classroom. We broke the ties at random in case of greater than three individuals among the \textit{top three}. Having made referrals among the \textit{top three} was the basis for earning a bonus in this part of the experiment.  

% \subsubsection{Baseline}
% Participants in the \textbf{baseline} treatment were asked to make three referrals for each skill measure. The order of the referrals for each skill measure were shown at random. To refer an individual, participants used a list where the names of their classmates were presented. They could tick the box in front of a name to refer that individual. As long as participants referred to three individuals per skill measure, they were free to pick anyone in their classroom. Participants could refer the same three individuals twice, but they were incentivized to choose the individuals they believed would be among the \textit{top three} for their classroom: 20\% of participants in each classroom were randomly selected to take part in a lottery to win a bonus worth 100,000 Pesos (about \$26). Among those chosen for the lottery, we randomly selected one skill measure and picked one of the participant's three referrals for that skill measure. If the selected referral was among the \textit{top three}, the participant received the bonus. 

% \subsubsection{Quota}
% Participants in the \textbf{quota} treatment were also asked to make three referrals for each skill measure. The order of the referrals for each skill measure were shown at random. The incentives for earning a bonus remained the same for the first two individuals among the \textit{top three}, i.e., participants referred individuals based on their performance in each skill measure regardless of their SES levels. However, the third member of the \textit{top three} for both skill measures would instead be the best performer among the lower SES participants (strata 1 to 3). We ranked each participant's correct answers for each skill measure from the highest to the lowest number, this time based on their SES: Either low-SES, consisting of strata 1,2,3 or high-SES, strata 4,5,6. We allocated one spot among the \textit{top three} to the best performer among the low SES. Once again, we broke ties at random. The remaining two slots of the \textit{top three} for that classroom were selected regardless of the SES. The purpose of this treatment was to introduce a quota to increase referrals of lower social classes.   

\subsection{Procedures}\label{Procedures}
% \textcolor{red}{1.4
% % \begin{itemize}
% %     \item - Explain that payoffs are calculated as three different bonuses. Bonus 1 is given to 10\% of students in each class for their performance in the tests. For this we randomly select n students in each class, for each we choose a test and their performance determines the probability of winning the bonus
% %     \item - Explain for the second bonus we chose 20\% of the class, for each the test is selected and then one of the referrals for that test. If the referral is in the top3 he earns the bonus
% %     \item - Explain for bonus 3 we do the social class [actually then you may want to explain this in the task part on top]
% % \end{itemize}
% % }

We collected data from 36 classrooms at Universidad UNAB for 20 minutes during an undergraduate-level course. The selected course is compulsory, and students take it at their convenience during three years of undergraduate studies. There were 36 asynchronous sessions dedicated for data collection. Our local partner, the Social Bee Lab, scheduled visits during the month of  April 2024 and recruited research assistants to administer the study protocol in each classroom. The entire data collection took place just before the exam period, allowing students more than four months to get to know each other. By the time of the data collection, the total number of enrolled students (i.e., eligible participants) stood at 864. Students' age varied between 17 to \textcolor{red}{22}. \textcolor{red}{X\%} of participants were female and \textcolor{red}{Y\%} belonged to the low-SES group (strata 1-3). Details of classroom compositions and descriptive statistics for the sample can be found in \textcolor{red}{Appendix 6}. The average expected payoff per participant was \$6, and it sat comfortably above the average hourly wage in Colombia for a 20 minute experiment.\footnote{The average monthly wage in Colombia is about \$1030, corresponding to an hourly wage of \$2. Calculations are made with an exchange rate of \$1 U.S. Dollar = 3700 Colombian Pesos.}

The incentive structure is unique for each part of the experiment. For the skill tests, 10\% of participants in each classroom (rounded upwards) are randomly selected to take part in a lottery to win a bonus worth 100,000 Pesos (about \$26). The chances to win the lottery weakly increases with the number of correct answers in each test. Among the participants selected for the lottery, we randomly select one skill test, and draw from the distribution determined by their answers to pay bonuses. In the 18-item Raven's test, every correct answer increases the chances to win by 5.5\% ($=1\backslash18$), and in the 36-item MRMET, every correct answer increases the chances to win by 2.75\% ($=1\backslash36$), conditional on that test being randomly selected for the lottery. 

In part two, incentivization works by referring high performing individuals within the classroom. With two individual skill measures available for each participant, we rank participants by the number of correct answers. For each skill measure in every classroom, we designate the \textit{top 3} performers. In case of ties within the \textit{top 3}, order is assigned randomly. 

For the \textbf{baseline} treatment, 20\% of participants in each classroom are randomly selected to take part in a lottery to win a bonus worth 100,000 Pesos (about \$26). Among the participants selected for the lottery, we randomly select one skill test, and draw one among three referrals. If the selected referral is in the \textit{top 3}, participants receive the bonus. In the case there are two top 3 referrals (one for each skill measure), participants win the lottery regardless of which skill test is selected. The chances to win weakly increases by one-third ($=33\%$) for each \textit{top 3} referral, conditional on that skill test being randomly selected for the lottery. 

For the \textbf{quota} treatment, we rank participants based on their SES and the number of correct answers in each test, using the individual stratum classification. The first position in the \textit{top 3} is reserved for the highest performing low-SES (strata 1 to 3) classmate, while the two remaining positions are chosen regardless of SES. In case the highest performing low-SES classmate is also the highest performing classmate, we take the second highest performer for that position. 20\% of participants in each classroom are randomly selected to take part in a lottery to win a bonus worth 100,000 Pesos (about \$26). Among the participants selected for the lottery, we randomly select one skill test, and draw one among three referrals. If the selected referral is among the \textit{top 3}, participants win the lottery. The chances to win weakly increases by one-third ($=33\%$) for each \textit{top 3} referral, conditional on that skill test being randomly selected for the lottery. However, there is an additional incentive to consider low-SES classmates because of the \textit{top 3} composition. It is now a dominant strategy to refer at least one high-performing low-SES regardless of the expected \textit{top 3} distribution in terms of SES. \textcolor{cyan}{Should I work on formalizing this?}

In part three, incentivization depended on the ability to recognize low-SES classmates independent of performance. 10\% of participants in each classroom are randomly selected to take part in a lottery to win a bonus worth 100,000 Pesos (about \$26). Among the participants selected for the lottery, we randomly draw one among three referrals. If the selected referral is from the low-SES group, participant wins the lottery. The chances to win strictly increases by one-third ($=33\%$) for each \textit{top 3} referral.

\textcolor{cyan}{Not sure if the incentives separated from the treatments belong here, I feel they make the reading more complicated. But having a dedicated space for the incentives allows me to write more in depth at the same time. What do you think?}

% \subsubsection{Timeline}
% The following sequence of events was identical for every classroom selected for the experiment.  The day before the research assistants' visit, every student in a given classroom received an email in their student mailbox. This email invited them to participate in the experiment, broadly described the study, and informed them about potential financial compensation.  Alongside the invitation, the email included a unique password specific to their classroom. This password served two purposes: It allowed students to access the correct experiment on the following day and prevented them from accessing experiments meant for other classrooms. On the next day, research assistants visited the classroom and projected a QR code for the experiment onto the whiteboard while reading the study description aloud. Students then used their mobile phones to scan the QR code and access the online experiment by typing in their unique classroom password. The experiment began with written instructions, followed by the informed consent form, skill measures, and referral tasks.

% \section{Results}


% \section{Limitations and Conclusion}


\bibliographystyle{apacite}
\bibliography{references}

\end{document}
